\section{Orações Diversas}

\subsection{Pro Summo Pontifice}
\begin{paracol}{2}\latim{
℣. Orémus pro Pontífice nostro {\redx N.}
}\switchcolumn\portugues{
℣. Oremos pelo nosso Pontífice {\redx N.}
}\switchcolumn*\latim{
℟. Dóminus consérvet eum, et vivíficet eum, et beátum fáciat eum in terra, et non tradat eum in ánimam inimicórum ejus.
}\switchcolumn\portugues{
℟. Que o Senhor o conserve e vivifique, que o faça santo na terra e o não entregue à vontade dos seus inimigos.
}\switchcolumn\portugues{
℟. Pater Noster... Ave Maria...
}\switchcolumn\portugues{
℟. Pai Nosso... Ave Maria...
}\switchcolumn\portugues{
℟. Deus, ómnium fidélium pastor et rector, fámulum tuum {\redx N.}, quem pastórem Ecclésiæ tuæ præésse voluísti, propítius réspice: da ei, quǽsumus, verbo et exémplo, quibus præest, profícere: ut ad vitam, una cum grege sibi crédito, pervéniat sempitérnam. Per Christum, Dóminum nostrum. Amen.
}\switchcolumn\portugues{
℟. Ó Deus, pastor e guia de todos os fieis, olhai misericordiosamente sobre vosso Servo {\redx N.} que escolheestes para pastor da vossa Igreja; concedei-lhe que, com sua palavra e exemplo, encaminhe o rebanho que lhe confiastes, e juntamente com ele, possa alcançar a vida eterna. Por Cristo Senhor nosso. Amen.
}\end{paracol}

\subsection{Prece a São Miguel Arcanjo}\label{precesaomiguel}
\begin{paracol}{2}\latim{
\rlettrine{P}{rinceps} gloriosíssime cæléstis militiæ, sancte Michaël Archangele, defénde nos in prælio advérsus príncipes et potestátes, advérsus mundi rectóres tenebrárum harum, contra spirituália nequitiæ, in cæléstibus.
}\switchcolumn\portugues{
\rlettrine{G}{loriosíssimo} Príncipe da Milícia Celeste, São Miguel Arcanjo, defendei-nos «no nosso combate contra os principados e potestades, contra os príncipes do mundo tenebroso, contra as hostes espirituais da iniquidade nas regiões celestes».
}\switchcolumn*\latim{
Veni in auxilium hominum; quos Deus ad imáginem similitúdinis suæ fecit, et a tyránnide diáboli emit prétio magno.
}\switchcolumn\portugues{
Vinde em auxílio dos homens, que Deus criou à sua imagem e semelhança e que remiu, por alto preço, da tirania do demónio.
}\switchcolumn*\latim{
Te custódem et patrónum sancta venerátur Ecclésia; tibi trádidit Dóminus ánimas redemptórum in supérna felicitáte locándas. Deprecáre Deum pacis, ut cónterat sátanam sub pédibus nostris, ne ultra váleat captivos tenére hómines, et Ecclésiæ nocére.
}\switchcolumn\portugues{
A Santa Igreja vos venera como Guarda e Patrono, a vós, Deus confiou as almas remidas destinadas a ter assento na suprema Felicidade. Rogai ao Deus da Paz que esmague Satanás debaixo dos nossos pés, que ele não possa mais reter os homens cativos e infligir males à Igreja.
}\switchcolumn*\latim{
Offer nostras preces in Conspéctu Altíssimi, ut cito anticipent nos misericórdiæ Dómini, et apprehéndas dracónem, serpéntem antiquum, qui est diábolus et sátanas, et ligátum mittas in abyssum, ut non sedúcat ámplius gentes.
}\switchcolumn\portugues{
Oferecei as nossas preces ao Altíssimo, a fim de que, sem demora, possa atrair a sua misericórdia sobre nós; prendei «o dragão, a antiga serpente que é o Diabo e Satanás», lançai-o acorrentado no abismo, «que ele mais não possa seduzir as nações».
}\end{paracol}

\subsection{Ao Anjo da Guarda}
\subsubsection{Ángele Dei}\label{angeledei}
\begin{paracol}{2}\latim{
\slettrine{Á}{ngele} Dei, qui custos es mei, me, tibi commíssum pietáte supérna, illúmina, custódi, rege et gubérna. Amen.
}\switchcolumn\portugues{
\rlettrine{S}{anto} Anjo do Senhor, meu zeloso guardador, pois que a ti me confiou a piedade divina, hoje e sempre me governa, rege, guarda e ilumina. Amen.
}\end{paracol}

\subsection{Fidelíssimo Companheiro}
\rlettrine{M}{eu} fidelíssimo companheiro, a quem Deus destinou para minha guarda, protector e defensor meu, muitas graças Vos dou por me haverdes livrado de tantos perigos; guiai-me no caminho do céu, preservai-me de todo o pecado e de toda desgraça: apresentai à divina Majestade as minhas orações e obras pias, os meus trabalhos e aflições, e fazei que em graça passe desta à vida eterna. Amen.

\subsection{Oração Universal do Papa Clemente XI}
\begin{paracol}{2}\latim{
\blettrine{C}{redo} Domine, sed credam firmius; spero, sed sperem securius; amo, sed amem ardentius; doleo, sed doleam vehementius.
}\switchcolumn\portugues{
\blettrine{S}{enhor,} creio em Vós, fazei que creia com mais firmeza; espero em Vós, fazei que espere com mais confiança; amo-Vos, aumentai o meu amor; arrependo-me, avivai a minha dor.
}\switchcolumn*\latim{
Adoro te ut primum principium; desidero ut finem ultimum; laudo ut benefactorem perpetuum; invoco ut defensorem propitium.
}\switchcolumn\portugues{
Adoro-Vos como primeiro princípio; desejo-Vos como último fim; exalto-Vos como benfeitor perpétuo; invoco-Vos como defensor propício.
}\switchcolumn*\latim{
Tua me sapientia dirige, iustitia contine, clementia solare, potentia protege.
}\switchcolumn\portugues{
Dirigi-me com vossa sabedoria; atai-me com vossa justiça; consolai-me com vossa clemência; protegei-me com vosso poder.
}\switchcolumn*\latim{
Offero tibi, Domine cogitanda, ut sint ad te; dicenda, ut sint de te; facienda, ut sint secundum te; ferenda, ut sint propter te.
}\switchcolumn\portugues{
Ofereço-Vos os meus pensamentos, para que se dirijam a Vós; minhas palavras, para que falem de Vós; minhas obras, para que sejam vossas; minhas contrariedades, para que as aceite por Vós.
}\switchcolumn*\latim{
Volo quidquid vis, volo quia vis, volo quomodo vis, volo quamdiu vis.
}\switchcolumn\portugues{
Quero o que quereis, quero porque o quereis, quero como o quereis, quero enquanto o queirais.
}\switchcolumn*\latim{
Oro, Domine, intellectum illumines, voluntatem inflammes, cor emundes, animam sanctifices.
}\switchcolumn\portugues{
Senhor, peço-Vos que ilumineis a minha mente, inflameis a minha vontade, limpeis o meu coração, santifiqueis a minha alma.
}\switchcolumn*\latim{
Defleam praeteritas iniquitates, repellam futuras tentationes, corrigam vitiosas propensiones, excolam idoneas virtutes.
}\switchcolumn\portugues{
Que me afaste das faltas passadas, rejeite as tentações futuras, corrija as más inclinações, pratique as virtudes necessárias.
}\switchcolumn*\latim{
Tribue mihi, bone Deus, amorem tui, odium mei, zelum proximi, contemptum mundi.
}\switchcolumn\portugues{
Concedei-me, Deus de bondade, amor por Vós, ódio por mim, zelo pelo próximo, desprezo pelo mundano.
}\switchcolumn*\latim{
Studeam superioribus oboedire, inferioribus subvenire, amicis consulere, inimicis parcere.
}\switchcolumn\portugues{
Que saiba obedecer aos superiores, ajudar os inferiores, acolher os amigos, perdoar os inimigos.
}\switchcolumn*\latim{
Vincam voluptatem austeritate, avaritiam largitate, iracundiam lenitate, tepiditatem fervore.
}\switchcolumn\portugues{
Que vença a sensualidade com a mortificação, a avareza com a generosidade, a ira com a bondade, a tibieza com a piedade.
}\switchcolumn*\latim{
Redde me prudentem in consiliis, constantem in periculis, patientem in adversis, humilem in prosperis.
}\switchcolumn\portugues{
Fazei-me prudente nos conselhos, constante nos perigos, paciente nas contrariedades, humilde na prosperidade.
}\switchcolumn*\latim{
Fac, Domine, ut sim in oratione attentus, in epulis sobrius, in munere sedulus, in proposito firmus.
}\switchcolumn\portugues{
Senhor, fazei-me atento na oração, sóbrio na comida, perseverante no trabalho, firme nos propósitos.
}\switchcolumn*\latim{
Curem habere innocentiam interiorem, modestiam exteriorem, conversationem exemplarem, vitam regularem.
}\switchcolumn\portugues{
Que procure ter inocência interior, modéstia exterior, conversa exemplar, vida ordenada.
}\switchcolumn*\latim{
Assidue invigilem naturae domandae, gratiae fovendae, legi servandae, saluti promerendae.
}\switchcolumn\portugues{
Que lute por dominar a minha natureza, fomentar a graça, servir a vossa lei e obter a salvação.
}\switchcolumn*\latim{
Discam a te quam tenue quod terrenum, quam grande quod divinum, quam breve quod temporaneum, quam durabile quod aeternum.
}\switchcolumn\portugues{
Que aprenda de Vós como é pouco o terreno, como é grande o divino, como é breve o tempo, como é duradouro o eterno.
}\switchcolumn*\latim{
Da mortem praeveniam, iudicium pertineam, infernum effugiam, paradisum obtineam.
}\switchcolumn\portugues{
Fazei-me preparar a morte, temer o juízo, evitar o inferno e alcançar o Paraíso.
}\switchcolumn*\latim{
Per Christum Dominum nostrum. Amen.
}\switchcolumn\portugues{
Por Cristo Nosso Senhor. Amen.
}\end{paracol}

\subsection{Bom Humor}
\rlettrine{C}{oncedei-me,} Senhor, uma boa digestão, mas também algo para digerir. Concedei-me um corpo saudável, e o bom humor necessário para o manter. Concedei-me uma alma simples que saiba valorizar o que é bom, que se não amedronte diante do mal, e que encontre os meios para voltar a colocar as coisas no seu lugar. Dai-me uma alma que não conheça o aborrecimento, nem as murmurações, nem os suspiros, nem os lamentos, nem preocupações excessivas com esse obstáculo chamado «eu». Concedei-me, Senhor, o dom do sentido de humor. Permiti-me a graça de aproveitar o riso para que saboreie nesta vida um pouco de alegria, e possa partilhá-la com os outros. Amen.

\subsection{Oração de São Francisco}
\rlettrine{S}{enhor,} fazei de mim um instrumento da vossa paz; onde houver ódio, que eu leve o amor; onde houver ofensa, que eu leve o perdão; onde houver discórdia, que eu leve a união; onde houver dúvida, que eu leve a fé; onde houver erro, que eu leve a verdade; onde houver desespero, que eu leve a esperança; onde houver tristeza, que eu leve a alegria; onde houver trevas, que eu leve a luz. Ó Mestre, fazei que eu procure mais consolar que ser consolado, compreender que ser compreendido, amar que ser amado. Pois é dando que se recebe, é perdoando que se é perdoado, e é morrendo que se vive para a vida eterna. Amen.

\subsection{Conversão dos Pecadores}
\begin{paracol}{2}\latim{
\rlettrine{Æ}{terne} rerum omnium effector Deus, memento abs te animas infidelium procreatas, easque ad imaginem et similitudinem tuam conditas. Memento Jesum, Filium tuum, pro illorum salute atrocissimam subiisse necem. Noli, quæso, Domine, ultra permittere, ut Filius tuus ab infidelibus contemnatur, sed precibus sanctorum virorum et Ecclesiæ, sanctissimi Filii tui Sponsæ, placatus, recordare misericordiæ tuæ et, oblitus idololatriæ et infidelitatis eorum, effice ut ipsi quoque agnoscant aliquando quem misisti Dominum Jesum Christum, qui est salus, vita et resurrectio nostra, per quem salvati et liberati sumus, cui sit gloria per infinita sæcula sæculorum. Amen.
}\switchcolumn\portugues{
\rlettrine{D}{eus} eterno, Criador de todas as coisas, lembrai-Vos que as almas dos infiéis são obras de vossas mãos, e que são feitas à vossa imagem e semelhança. Vede, porém, Senhor, como em desonra do vosso Nome o inferno se enche destas almas. Lembrai-Vos que Jesus Cristo, vosso Filho, derramou todo seu Sangue e padeceu uma morte atrocíssima por elas. Não permitais, pois, Senhor, que o vosso Filho seja por mais tempo desprezado pelos infiéis. Deixai-Vos antes mover à piedade pelas orações de vossos Santos e da Igreja, esposa de vosso Santíssimo Filho. Lembrai-Vos da vossa misericórdia e, esquecendo a sua idolatria e infelicidade, fazei que também eles conheçam Jesus Cristo, Nosso Senhor, que é nossa Salvação, Vida e Ressurreição, pelo qual fomos livres e salvos do inferno, e a quem seja dado honra, glória e louvor para sempre. Amen.
}\end{paracol}

\subsection{Jaculatórias}
Meu Jesus, misericórdia!\par
Jesus, manso e humilde de coração, fazei o meu coração semelhante ao vosso.\par
Doce Coração de meu Jesus, fazei que eu Vos ame cada vez mais.\par
Seja feita, louvada e eternamente exaltada a justíssima, altíssima e amabilíssima vontade de Deus em todas as coisas.\par
Ó Maria concebida, sem pecado. Rogai por nós que recorremos a Vós.\par
Ó meu Jesus, perdoai-nos, livrai-nos do fogo do inferno. Levai as almas
todas para o Céu e socorrei principalmente as que mais precisarem\par
Ó Jesus, manso e humilde de coração, fazei o nosso coração semelhante
ao vosso.\par
Meu Deus, eu creio, adoro, espero e amo-Vos. Peço-Vos perdão para os que
não crêem, não adoram, não esperam e não Vos amam.\par
Senhor, valei-me.\par
Senhor, seja feita a vossa santíssima vontade.\par
Maria, Auxílio dos Cristãos, rogai por nós.\par
Meu Jesus, eu quero ser todo vosso.\par
Meu Jesus, misericórdia.\par
Doce Coração de Jesus, que tanto nos amais, fazei com que eu Vos ame cada vez mais.\par
Doce Coração de Maria, sede minha salvação.\par
Santa Bárbara Bendita, que no céu está escrita com papel e água benta, nos livre desta tormenta.\par
São Miguel Arcanjo, defendei-nos no combate.\par

\subsection{Cruz de S. Bento}
\begin{nscenter}
\includegraphics[width=.8\textwidth]{media/bento}
\end{nscenter}
\begin{paracol}{2}\latim{
{\redx C.S.P.B.:} Crux Sancti Patris Benedicti.
}\switchcolumn\portugues{
{\redx C.S.P.B.:} Cruz do Santo Patriarca Bento.
}\switchcolumn*\latim{
{\redx C.S.S.M.L.:} Crux Sancta Sit Mihi Lux.
}\switchcolumn\portugues{
{\redx C.S.S.M.L.:} A Cruz Santa seja a minha Luz.
}\switchcolumn*\latim{
{\redx N.D.S.M.D.:} Non Draco Sit Mihi Dux.
}\switchcolumn\portugues{
{\redx N.D.S.M.D.:} Que o Dragão não seja meu Senhor.
}\switchcolumn*\latim{
{\redx V.R.S.:} Vade Retro Satana!
}\switchcolumn\portugues{
{\redx V.R.S.:} Retira-te, Satanás!
}\switchcolumn*\latim{
{\redx N.S.M.℣.} Numquam Suades Mihi Vana!
}\switchcolumn\portugues{
{\redx N.S.M.℣.} Não me aconselhes loucuras!
}\switchcolumn*\latim{
{\redx S.M.Q.L.:} Sunt Mala Quæ Libas.
}\switchcolumn\portugues{
{\redx S.M.Q.L.:} São maldades o que me apresentas.
}\switchcolumn*\latim{
{\redx I.V.B.:} Ipse Venena Bibas.
}\switchcolumn\portugues{
{\redx I.V.B.:} Tu mesmo bebe esses venenos.
}\switchcolumn*\latim{
Amen.
}\switchcolumn\portugues{
Amen.
}\end{paracol}

\subsection{Renovação das Promessas do Baptismo}
\rlettrine{R}{enuncio} hoje a Satanás, a todas suas obras e a todas suas pompas. A Vós, meu Jesus, prometo fé e fidelidade, quero ficar fiel ao vosso amor até à morte. Afasta-te de mim, Satanás! Pertenço Jesus, e só a Jesus quero servir. Meu amabilíssimo Jesus, não permitais que hoje por algum pecado me separe de Vós.

\subsection{Requiem Æternam}\label{requiem}
\begin{paracol}{2}\latim{
\rlettrine{R}{équiem}ætérnam dona eis, Dómine, et lux perpétua lúceat eis. Requiéscant in pace. Amen.
}\switchcolumn\portugues{
\rlettrine{D}{ai-lhes,} Senhor, o eterno descanso entre os esplendores da luz perpétua. Descansem em paz. Amen.
}\end{paracol}

\subsection{Aceitação da Morte}
\rlettrine{M}{eu} Deus e meu Pai, Senhor da vida e da morte, que, para justo castigo das nossas culpas, com um decreto imutável determinastes que todos os homens haviam de morrer, olhai para mim, aqui prostrado diante de Vós. Detesto com todo o coração as minhas culpas passadas, pelas quais mereci mil vezes a morte, que aceito agora com o fim de expiá-las e para obedecer à vossa amável vontade. Gostosamente morrerei, Senhor, no momento, no lugar e do modo que Vós quiserdes, e aproveitarei até essa altura os dias que me restam de vida para lutar contra os meus defeitos e aumentar o meu amor por Vós, para quebrar os laços que atam o meu coração às criaturas, para preparar a minha alma para comparecer na vossa presença; e desde agora me abandono sem reservas nos braços da vossa paternal Providência.

\subsection{Obter uma Boa Morte}
\rlettrine{M}{eu} Criador e meu Pai, peço-Vos a mais importante das vossas graças: a perseverança final e uma morte santa. Por maior que tenha sido o abuso da vida que me destes, concedei-me a graça de vivê-la desde agora e de terminá-la no vosso santo amor. Que eu morra como os Santos Patriarcas, abandonando sem tristeza este vale de lágrimas, para ir gozar o descanso eterno na minha verdadeira pátria. Que eu morra como o glorioso S. José, acompanhado por Jesus e Maria, pronunciando estes nomes dulcíssimos, que espero bendizer por toda a eternidade. Que eu morra como a Virgem Imaculada, dentro da mais pura caridade e com o desejo de me unir ao único objecto dos meus amores. Que eu morra como Jesus na Cruz, plenamente identificado com a vontade do Pai, feito um holocausto por amor. Jesus, morto por mim, concedei-me a graça de morrer com um acto de perfeita caridade para convosco. Santa Maria, Mãe de Deus, rogai por mim agora e na hora da minha morte. S. José, meu Pai e Senhor, alcançai-me a graça de morrer com a morte dos justos.

\subsection{Momento da Morte}
\rlettrine{S}{enhor,} meu Deus, de boa vontade aceito desde já, como vinda das vossas mãos, qualquer espécie de morte que quiserdes enviar-me, com todas suas angústias, penas e dores.\par
℣. Amado Jesus, José e Maria,\par
℟. Dou-Vos o coração e a alma minha.\par
℣. Amado Jesus, José e Maria,\par
℟. Assisti-me na última agonia.\par
℣. Amado Jesus, José e Maria,\par
℟. Expire em paz entre Vós a alma minha.\par
